\documentclass[../main.tex]{subfiles}
\begin{document}
%==========================================TIÊU ĐỀ TẬP========================================
\begin{table}[h]
  \begin{adjustwidth}{-1cm}{}
    \begin{tabular}{>{\centering\arraybackslash}p{6cm}|>{\raggedright\arraybackslash}p{12.5cm}}
      \multirow{5}{6cm}
      % Cột bên trái: ÉPISODE và số tập
      {\\[-40pt]
        \flaregothic\fontsize{20pt}{20pt}\selectfont \centering
        {\contour{errortwo}{\textcolor{black}{ÉPISODE}}}\color{black}\\
        \burbank\fontsize{72pt}{72pt}\selectfont
        \includegraphics[height=3cm]{../../../../Image/Book?/number/num5.png}
        \\[18pt]
        % \flaregothic\fontsize{14pt}{14pt}\selectfont
        % \color{epattr}BIENVENUE, \uline{\color{Banpire} Mel} !
        % \flaregothic\fontsize{18pt}{18pt}\selectfont
        % \color{epattr}ÉPISODE PERDU
        \color{black}
      }
       &
      % Dòng thứ nhất: tên tập tiếng Nhật
      {\fotskip\fontsize{18pt}{18pt}\selectfont \ruby{美}{み}\ruby{空}{そら}\ruby{時}{し}\ruby{乃}{の}の\ruby{呪}{のろ}い
      } \\[6pt]
      % Dòng thứ hai: tên tập tiếng Anh
       & {\cafeteria\fontsize{18pt}{18pt}\selectfont Misora Shino's Curse}                                 \\[4pt]
      % Dòng thứ ba: tên tập tiếng Việt
       & {\vnmsans\fontsize{18pt}{18pt}\selectfont Lời ngiuyền của Misora Shino}                           \\[8pt]
      % Dòng thứ tư: Nhân vật xuất hiện
       & {\caviar{\fontsize{14pt}{14pt}\selectfont Track used: The Space}}
      \\[6pt]
      % Dòng cuối cùng: Thời gian diễn ra của tập (thường sẽ là ngày phát hành)
       & {\continuum\fontsize{16pt}{16pt}\selectfont Ngày phát hành: 26/03/2022}                           \\[8pt]
    \end{tabular}
  \end{adjustwidth}
\end{table}
%===========================================NỘI DUNG==========================================
\color{black}

\pov{Haragen Kyuen}

``Còn điều gì níu giữ cậu nữa? Mau cùng về thôi, \emph{những học sinh mới}.''

Giọng Shino vang vọng trong đầu tôi, từng chữ như một lưỡi dao lạnh lẽo rạch qua suy nghĩ. Cô ta đứng ngay trước mặt, hai tay chắp ra sau, đôi mắt rỗng không như đã đánh mất hết thảy sự sống, nụ cười ngờ nghệch cố định trên khuôn mặt. Tôi không thể đọc nổi cô đang nghĩ gì, hay có còn người trong cái vỏ bọc ấy nữa không.

\img{e1-5a}

``Rốt cuộc\dots\ cô đang làm cái gì ở đây vậy?''

\dots

\dots

Không có câu trả lời.

Và rồi, vẫn giữ nụ cười ma mị kia, Shino quay lưng bỏ đi, hướng về cánh cửa trượt đang chờ ở cuối hành lang. Bóng lưng ấy cứ xa dần, bước chân đều đặn một cách khó chịu, như thể đã quen thuộc con đường này từ lâu, còn tôi chỉ là kẻ lạc lõng bị kéo theo vào một trò chơi chẳng hiểu đầu đuôi.

``Về thôi''? Cô ta nói như thể nơi này vốn dĩ đã định sẵn là nhà của chúng tôi. Nhưng tại sao? Tại sao Shino lại xuất hiện ngay lúc ngôi trường sụp đổ? Tại sao cô ta nhìn tôi với ánh mắt chết chóc đó?

Và\dots\ Shino này\dots\ giống như những gì đã xảy ra trong giấc mơ đêm đó. Nhưng cô ta là ai? Cô ta có sự liên kết gì với ngôi trường ấy? Tại sao mọi chuyện lại thành ra như thế này?

Những câu hỏi cứ xoắn lấy nhau trong đầu, để rồi khi tôi nhận ra, bản thân đã đứng chết lặng, chỉ còn lại bóng lưng xa dần của Shino và một cánh cửa khép kín trước mặt.

Tôi không biết\dots\ điều gì thực sự đang xảy ra.

\dots

\dots

``\dots-chan!''

``\dots-kun!''

Giọng của hai người khác.

Tôi nghe tiếng gọi vọng lại giữa cơn mơ hồ của bản thân. Ban đầu chỉ là những âm thanh rời rạc, văng vẳng như vọng từ xa xăm.

Nhưng rồi, một tiếng quen thuộc vang lên rõ rệt. Hai giọng nói này rất quen. Không phải từ phe địch.

``\textbf{Nii-chan!}''

Ngay sau đó, một giọng khác nối tiếp, gấp gáp không kém.

``\textbf{Kyuen-kun!}''

Mi mắt tôi nặng trĩu, nhưng rồi cũng dần mở ra. Thứ đầu tiên đập vào mắt tôi không phải những ảo ảnh trong mơ, mà là gương mặt lo lắng của cô em gái ngốc nghếch của tôi. Ngay cạnh đó là một cô gái khác - Sora-san.

Tôi thoáng giật mình, tim khựng lại một nhịp. Bởi trong khoảnh khắc ấy, tôi tưởng mình thấy Shino. Mái tóc, khuôn mặt, cái dáng người mảnh khảnh kia\dots\ quá giống. Chỉ đến khi ánh mắt nâu của cô gái chạm vào tôi, khác hẳn với đôi mắt lạnh lẽo của Shino, tôi mới nhận ra mình đã nhầm.

``Nii-chan\dots\ anh tỉnh rồi,'' Saikyou run run nói, rồi vội nhìn khắp thân thể tôi. Em khẽ thở phào khi thấy tôi chỉ bị xây xát ngoài da.

Tôi gượng ngồi dậy, khẽ hỏi lại: ``Còn em\dots\ và Sora-san thì sao? Có bị thương gì không?''

Em lắc đầu, nở một nụ cười nhẹ nhõm. Quần áo em dính đầy bụi, da bị xướt xát nhẹ, nhưng may mắn không có gì nghiêm trọng.

Sora thì chỉ mất chiếc áo khoác ngoài, giờ chỉ còn lại bộ sơ-mi trắng ngắn tay và chiếc váy ca-rô xanh, giống như bộ mà em gái tôi đang mặc.

Cả ba chúng tôi im lặng một thoáng, để hơi thở lắng xuống, rồi mới chậm rãi quan sát xung quanh.

Và rồi, tôi sững người.

Thứ hiện ra trước mắt không còn là hành lang vừa chạy qua, hay bất cứ ở lớp học 1-C đó.

Mà lại là một nơi\dots\ không khác với thứ mà tôi đã trải qua trong đêm hôm đó.

\img{e1-5b}

Từ tường, sàn cho đến trần nhà, tất cả mọi thứ đều được bao phủ bởi những mẩu gỗ đã mục dần theo thời gian. Những bóng đèn huỳnh quang cũ kỹ phát ra thứ ánh sáng yếu ớt, chập chờn như ngọn nến trước gió. Dọc theo hành lang là hai dãy phòng học, mỗi phòng được ngăn cách bởi những cánh cửa trượt kiểu truyền thống với những ô lưới nhỏ đã ngả màu thời gian.

Phía trên mỗi lớp học, những tấm biển tên đã bạc phếch, chỉ còn sót lại vài vệt mực đen mờ ảo khó đọc. Trên những cánh cửa, các ô kính vỡ nát để lại những khung gỗ nhô ra sắc nhọn, như thể đó là hàm răng của một sinh vật vô hình đang rình rập, chờ đợi con mồi tiếp theo bước qua ngưỡng cửa của nó.

Tôi nuốt khan, cảm giác lạnh sống lưng chạy dọc suốt từ gáy xuống chân.

Toàn bộ không gian này\dots\ đậm đặc mùi của thời đại đã qua, thứ hoài cổ u ám như thể lạc bước vào một ngôi trường từ thế kỷ trước. Hoàn toàn khác với Aogami mà chúng tôi biết.

Một sự trùng hợp tới khó tin so với những gì mà hai anh em tôi đã trải qua trong giấc mơ đêm đó.

Tôi nuốt khan, cảm thấy không khí xung quanh như đông đặc lại. Cả hành lang phủ một màu tối ẩm mốc, chỉ còn ánh sáng nhỏ bé từ điện thoại của ba chúng tôi lách ra từng vệt mỏng. Mùi gỗ ẩm mốc xen lẫn mùi sắt rỉ len lỏi trong hơi thở, khiến tôi chỉ muốn quay lưng chạy đi ngay lập tức.

Saikyou bước chậm rãi hơn thường ngày, đôi mắt cô bé ấy nhìn chăm chú vào từng mảng tường bong tróc. Tôi biết em ấy đang nhớ lại giấc mơ đêm trước khi nhập học, nơi mà cả hai anh em tôi đều đã trải qua. Giấc mơ về những dãy hành lang vỡ vụn, lớp học méo mó như bị kéo giãn, bàn ghế sắp xếp lộn xộn. Tôi run lên khi nhận ra khung cảnh này chẳng khác gì bản sao từ cơn ác mộng đó.

``Nii-chan\dots'' Saikyou khẽ gọi, giọng cô trầm xuống. ``Anh cũng nhìn thấy nơi này rồi\dots\ đúng không?''

Tôi im lặng, chỉ gật đầu. Câu trả lời ấy đủ để cả hai hiểu rằng, đây không phải sự trùng hợp đơn thuần. Có thể hành lang này trông tươm tất và sạch sẽ hơn so với cái mà chúng tôi thấy trong mơ, nhưng sự u ám này\dots\ không thể nhầm đi đâu được.

Sora lại có phản ứng khác hẳn so với hai anh em chúng tôi. Cô ấy không giấu nổi sự căng thẳng; bàn tay cầm điện thoại khẽ run, ánh sáng hắt lên khuôn mặt nghiêm nghị, đôi mắt không rời khỏi những ô cửa lớp học chập chờn đèn hai bên.

``Cẩn thận. Chúng ta không nên ở đây lâu đâu,'' giọng Sora hạ thấp, như thể chính cô đã từng bước đi trong hành lang này từ trước rồi.

Chúng tôi tiến về phía trước. Những căn lớp học hai bên đều sáng đèn, nhưng chẳng có lấy một bóng người. Bàn ghế bên trong ngổn ngang, mỗi cái quay một hướng, giống như có ai đó đã cố tình phá vỡ trật tự vốn có của nó. Âm thanh bước chân vang lên lạnh lẽo, vọng đi mãi trong khoảng không im lìm, khiến tôi rùng mình.

Đến cuối hành lang, trước mặt chúng tôi là một cánh cửa gỗ cũ kỹ. Saikyou thử dùng sức đẩy mạnh vài lần, nhưng cánh cửa chỉ rung lên rồi đứng im, kẹt chặt như bị khóa từ bên kia. Cô nhíu mày, lùi lại, hơi thở nặng nề.

``Chặt quá\dots\ Em không mở được\dots''

``Có vẻ như là có ai đó khóa trái rồi,'' Sora nói, rôi chiếu đèn pin về phía bên cạnh.

Một hành lang khác hiện ra. Trái ngược hoàn toàn, nơi ấy sáng trưng bởi ánh đèn huỳnh quang trắng nhợt. Chúng tôi ngập ngừng nhìn nhau, ánh sáng rõ ràng an toàn hơn bóng tối, nhưng trong ánh sáng ấy lại có điều gì khiến tôi khó chịu, như thể nó giả tạo, gượng gạo, giống một nụ cười đông cứng không chạm tới đáy mắt.

``Đi tiếp thôi nào, hai người,'' Sora nói, nhưng tôi nghe trong giọng cô không chỉ có sự cảnh giác cao độ dù đã quen đường đi nước bước, mà còn là\dots\ một nỗi sợ thầm kín nào đó mà cô đang cố nén lại.

Chúng tôi bước vào hành lang sáng. Mỗi bước đi đều nặng như kéo lê cả thân thể. Tiếng đèn huỳnh quang kêu rè rè, ánh sáng nhấp nháy loạn nhịp, soi bóng chúng tôi méo mó trên tường.

Cuối cùng, khi đi hết dãy hành lang ấy, trước mặt lại hiện ra một cánh cửa trượt gỗ khác. Chúng tôi dừng lại, lặng im nhìn nhau, rồi tôi hít một hơi thật sâu, bàn tay run run đặt lên tay nắm gỗ đã mòn vẹt.

``Cẩn thận đấy, Kyuen-kun,'' cô bạn tóc nâu lên tiếng nhắc nhở. Tôi gật đầu nhẹ ra dấu.

Tôi kéo sang một bên, và *XOẠCH*, âm thanh khô khốc vang vọng cả hành lang.

Phía sau cánh cửa lại là một dãy hành lang khác, sáng đèn, nhưng ánh sáng ấy chẳng hề khiến nơi này bớt đáng sợ. Mấy bóng huỳnh quang treo lủng lẳng trên trần bắt đầu chập chờn, lúc sáng lúc tối.

Hai bên lối đi, bàn ghế ngổn ngang, chất chồng lên nhau, có cái thì đổ ngang, có cái bị kẹt sát tường. Bên trong lớp, một sự trống trải bâu lấy, bàn ghế vẫn kê xộc xệch như thể nơi này đã bỏ hoang từ khi nào.

``Tại sao chỗ này lại bừa bộn vậy\dots'' cả hai anh em chúng tôi đồng thanh, tần ngần nhìn hành lang trông như thể vừa kết thúc một hội chợ trường mà chưa kịp đi dọn.

``Tớ không biết tại sao,'' Sora cất lời, ánh đèn pin rọi về phía trước, khuôn mặt lộ rõ vài giọt mồ hôi, ``nhưng chúng ta vẫn phải cảnh giác cao độ. Rất có thể \emph{cô ta} lại đang định giở trò gì nữa đây.''

Chúng tôi chỉ có thể ``ừm'' một tiếng nhẹ, thận trọng tiến về phía trước.

Chúng tôi len lỏi qua khoảng trống hẹp giữa đống gỗ mục, mỗi bước đi càng khiến bầu không khí vốn đã ngột ngạt thêm phần bí bách. Tôi vẫn có thể nghe thấy tiếng gió thổi nhẹ bên ngoài, càng khiến cho không khí bên trong đây rùng rợn hơn.

*ROẸT\dots*

Bất ngờ, một chiếc ghế ngay sát tường rít lên khe khẽ rồi tự mình trượt sang bên cạnh, như có ai vô hình vừa đẩy.

``Á\dots!'' Saikyou thót lên, nắm chặt lấy tay áo tôi. Tôi cũng sững cả người, tim đập thình thịch trong lồng ngực.

``Sao lại\dots'' tôi lắp bắp, không tin nổi vào mắt mình. Trong giấc mơ, tôi từng thấy những hình ảnh quái đản, nhưng tận mắt chứng kiến nó xảy ra ở đây\dots\ lại là chuyện khác hoàn toàn.

Sora thì vẫn đứng đó, mắt cô không rời khỏi chiếc ghế vừa nhúc nhích, ánh nhìn bình thản đến khó hiểu. ``Đừng hoảng loạn. Đây không phải lần đầu tôi thấy thứ này đâu,'' giọng cô nàng lạnh và chắc, trái ngược với sự hoang mang đang dâng trào trong chúng tôi.

Chúng tôi chưa kịp định thần thì một chiếc ghế khác bất ngờ lạch cạch chuyển động ngay trước mặt. Tôi giật mình lùi lại, nhưng đúng lúc đó một chiếc ghế khác bên cạnh cũng xoay ngang, chắn ngay lối đi. Tôi vấp phải chân ghế, mất thăng bằng.

``\textbf{Nii-chan!}'' em gái tôi kêu to, lao tới đỡ lấy tôi, cùng với Sora kịp thời giữ vai tôi kéo lại. Tôi thở dốc, tim như muốn nhảy khỏi lồng ngực.

``Nguy hiểm quá\dots\ chuyện này không có trong mơ đâu, Nii-chan\dots'' Saikyou khẽ run, nhưng vẫn nhìn tôi chờ lời khẳng định.

Tôi gượng gạo đáp, giọng lạc đi: ``\dots Không giống đâu. Trong mơ\dots\ chúng ta chỉ thấy cảnh tượng, thậm chí ta có thể biết được sơ qua những gì đã xảy ra nơi đây qua những mẩu giấy. Nhưng ở đây\dots\ mọi thứ đang thực sự chuyển động. Nhất là sau sự kiện sập trường mà ta đã trải nghiệm.''

Em tôi cúi đầu, thì thầm: ``Vậy\dots\ những thứ mà ta khám phá được, có khi nào\dots\ đều liên quan đến nơi này không?''

``Hừm\dots''

Tôi im lặng, trong đầu chợt lóe lên.

``\dots Cũng có thể lắm. Thậm chí anh có thể nghĩ\dots\ đây là ngôi trường Aogami từ thế kỷ trước.''

``Thế ạ?''

``Khả năng cao là vậy. Có thể sẽ có vài điểm sai khác, nhưng anh có thể khẳng định là như vậy.''

Mắt tôi bỗng bắt gặp một ánh sáng mờ phát ra từ phía trên một chiếc bàn lật nghiêng. Tôi bước lại gần, gạt mấy mảnh gỗ ra, và nhìn thấy một vật hình tròn làm bằng nhựa, với mặt kính trong suốt, trông như kính cường lực bảo vệ. Bên trên là một vòng kim loại bao quanh, có thể xoay được cả hai chiều, phát ra những tiếng  ``lạch tạch'' nhỏ nhẹ. Trên vành khung còn lồi ra vài nút nhỏ như để nhấn.

Ánh sáng chạp chờn của hành lang chiếu vào mặt kính bảo vệ mờ đục vì bụi ấy, dù dịu và yếu ớt, nhưng khiến tim tôi đập rộn lên.

Tôi lấy từ trong túi ra những thứ đã cất giữ từ trước (mà tôi vẫn ngạc nhiên là nó vẫn còn ở đây): mặt sau bằng thép không gỉ đen nhám với hình vuông đen kẹp giữa hai hình chữ nhật trắng có khắc ``Geat Frontier'' mà tôi nhặt trong giấc mơ đêm đó, và đoạn dây cao su màu đen nhám với các lỗ tròn.

``Hình như em cũng có dây gì đó đấy, Nii-chan\dots'' em tôi cất tiếng nhẹ.

``Hửm?''

``Hôm mà em với Touka-san và Inari-san về từ hôm trước đấy ạ\dots\ Em có nhặt được thêm một dây khác nữa.''

Em ấy lôi ra từ túi váy một cái dây đeo khác với phần đuôi có khung kim loại nhỏ cài một cái kim mảnh. ``Đây, Nii-chan. Em nghĩ là nó có thể là một phần của chiếc đồng hồ đó đấy\dots''

Tôi thử ghép chúng lại với nhau, cẩn thận luồn hai dây đeo vào hai cuộn kim loại đính ở mặt sau và mặt trước, rồi chỉnh lại một chút\dots

*CẠCH*

Vừa khít. Có vẻ như là\dots\ chúng sinh ra là để ăn khớp với nhau. Trước mắt tôi bây giờ là hình hài hoàn chỉnh của một chiếc đồng hồ đeo tay lạ lẫm - mà hình như tôi đã thấy nó ở đâu rồi, hoặc đúng hơn, đã từng sở hữu nó từ lúc nào.

Dù mặt kính cường lực chỉ phản chiếu bóng mờ của mặt tôi, không hề hiển thị gì, nhưng linh cảm mách bảo tôi rằng nó ẩn giấu một manh mối quan trọng. Có thể là về chính nơi này, hoặc thứ gì còn sâu hơn nữa.

Em gái tôi ngồi xổm xuống, nhìn chằm chằm chiếc đồng hồ rồi thì thầm: ``Thứ này\dots\ quan trọng lắm nhỉ, Nii-chan?''

Tôi gật đầu, nhưng không trả lời thêm. Bởi tôi cũng chẳng chắc.

Sora chỉ đứng một bên, quan sát không nói gì, nhưng ánh mắt cô ấy như đã hiểu rõ hơn chúng tôi. Khi thấy tôi cất những bộ phận của đồng hồ đeo tay lại vào túi, cô chỉ khẽ nhắc nhở: ``Ta cùng đi tiếp thôi.''

Chúng tôi lại tiếp tục bước đi, tránh qua những chiếc bàn ghế chặn đường. Hành lang dài bất tận, ánh sáng huỳnh quang mờ nhòe. Cuối cùng, một cánh cửa trượt gỗ khác hiện ra trước mắt.

``Đ-Để em mở cánh cửa này\dots'' Saikyou rụt rè tiến lên trước, thử dùng cả hai tay kéo nó sang bên. Nhưng\dots

*KÉT\dots*

\dots cánh cửa không nhúc nhích, dường như bị dính chặt vào khung. Cô thử thêm một lần nữa, toàn thân nghiêng hẳn về phía trước, nhưng vẫn vô ích.

``Tệ thật\dots'' tôi tặc lưỡi, quay sang Sora. ``Có lẽ chúng ta nên quay lại tìm lối khác thôi.''

Sora đưa tay chạm nhẹ lên bề mặt cánh cửa gỗ, ánh mắt như đang dò xét kỹ lưỡng, rồi khẽ lắc đầu: ``Không chắc đâu. Nhưng đúng là\dots\ nó bị khóa từ bên trong. Nếu quay lại, chưa chắc sẽ tìm được lối khác tốt hơn.''

Tôi thoáng nhìn về phía cô em gái. Im bặt, mặt tối sầm xuống, hai tay nắm chặt, trông như có vẻ uất ức lắm.

``Tránh ra.'' Em bất ngờ cất giọng, ngắn gọn nhưng đầy kiên quyết. ``Để em.''

Tôi và Sora đồng loạt nhìn cô bé, có linh cảm không hay. ``Khoan đã, Saikyou, em định--'' tôi chưa kịp nói hết câu.

Cô em gái tôi đã lùi lại vài bước, hít một hơi thật sâu, rồi bất thình lình lao lên.

``Này, khoan--!''

\textbf{*RẦM!*}

Tiếng đạp khô khốc vang dội cả hành lang. Cánh cửa gỗ vốn đã mục nát chẳng chịu nổi cú va chạm mạnh mẽ ấy, bật tung bản lề rồi ngã sập xuống nền. Trước mặt chúng tôi, một khoảng tối sâu thẳm mở ra, đen kịt như miệng hang nuốt chửng mọi ánh sáng.

Hai chúng tôi cùng toát mồ hôi hột, ngẩn người nhìn Saikyou đang thở hồng hộc. ``Em\dots\ em làm hơi quá rồi đấy\dots'' tôi nhăn mặt, tim còn chưa kịp bình ổn.

Đáp lại, em ấy chỉ khẽ nhún vai, che giấu sự run rẩy sau nụ cười gượng gạo: ``Đi tiếp thôi nào, hai người.''

\ct

Bên trong hành lang ấy, bóng tối đặc quánh đến mức ánh sáng từ điện thoại của chúng tôi chỉ soi được vài bước chân phía trước, xa hơn thì hoàn toàn bị nuốt chửng. Không gian tĩnh lặng, lạnh lẽo, chỉ còn tiếng bước chân vang vọng khẽ khàng.

``\dots Nii-chan, em không thích chỗ này chút nào.'' Saikyou lí nhí, mắt dán chặt vào vệt sáng yếu ớt.

``Anh cũng vậy. Cứ bám sát nhau, đừng tách ra.'' tôi dặn, tay nắm chắc cổ tay em gái.

Sora đi sau, giọng bình thản đến lạ: ``Hai người mới là người chưa quen thôi. Còn tớ thì\dots''

Cô dừng lại nửa nhịp, như đang tự cân nhắc. ``\dots Đã từng đi qua những hành lang giống thế này. Nhiều lần rồi.''

Tôi ngạc nhiên liếc nhìn: ``Ý cậu là\dots\ đã từng tới nơi này?''

``\dots Không hẳn, nhưng\dots\ đủ để biết. Những chỗ kiểu này hiếm khi chỉ đơn thuần là `trống rỗng'. Nơi này còn rình rập những điều kỳ quái hơn thế,'' cô đáp như đang gợi mở một cái gì đó khác, tay khẽ siết chặt chiếc điện thoại, ánh sáng rung rung.

Em tôichau mày, cố cứng cáp đáp lại: ``Vậy thì càng phải nhanh ra khỏi đây.''

``Tớ nghĩ hai cậu không nên ở đây. Để tớ giúp hai cậu tìm được đường khỏi nơi này.''

Chúng tôi bước từng bước nhỏ, cẩn trọng như đi trên băng mỏng. Bất chợt\dots

``Hi hi hi\dots\ Ha ha ha\dots''

Âm thanh mơ hồ vang lên từ đâu đó trong bóng tối. Tôi và Saikyou lập tức dựng tóc gáy, toàn thân như cứng đờ lại.

Thứ âm thanh này\dots\ tôi đã không nghe không dưới ba lần ở bên trong giấc mơ đêm đó. Những tiếng cười khúc khích của các học sinh đã từng theo học tại ngôi trường này.

``N-Nghe thấy gì không, Nii-chan\dots?'' Saikyou run giọng, bàn tay siết chặt lấy tay áo tôi.

``\dots Có, tiếng đó\dots\ anh biết,'' tôi nuốt khan, mắt căng ra tìm kiếm trong khoảng tối đặc quánh.

Sora lúc này cũng dừng lại, nhưng không hốt hoảng. Dẫu vậy, tôi thấy những giọt mồ hôi lăn dài trên thái dương cô. ``Tiếng này\dots\ không phải lần đầu tôi nghe.'' cô khẽ nói. ``Mọi người cẩn thận.''

Tiếng cười lại vang lên, lần này rõ ràng hơn, gần hơn. Khúc khích. Khúc khích. Những tiếng ``hi hi'', ``ha ha'' tưởng chừng vô hại nhưng lại khiến chúng tôi ám ảnh đến tận bây giờ.

Chúng tôi rụt rè bước tiếp, và rồi một căn phòng học sáng đèn bất ngờ hiện ra bên hành lang.

Ngay dưới cái đèn huỳnh quang chập chờn, một lớp học ở phía bên trái bất chợt sáng đèn. Tôi thấy có bóng người bên trong. Không phải chỉ một, hay hai\dots\ mà cả một lớp học đủ chỗ ngồi. Hai cánh cửa kéo bị đóng chặt, chẳng thể nào mở ra, như thể vừa bị hóa thành bức tường.

\img{e1-5c}

Tôi ngước về những người đang ngồi trong lớp học ấy. Tất cả mọi người\dots\ họ đều ngồi ngay ngắn, quay mặt về phía bảng đen. Không một tiếng nói, không một cử động. Không có giáo viên. Không có lý do gì để họ ngồi đó, trong căn phòng chết lặng này, chưa kể rằng bây giờ còn đang là ban đêm (tôi cho là vậy) nữa.

\dots Khoan. Tôi thấy cảnh tượng này rất quen thuộc. Trong cơn ác mộng đêm đó, tôi và Saikyou đã chạm trán với cái lớp học kỳ quặc kia.

Họ\dots\ quả nhiên không còn là ``người'' nữa.

Tôi đứng sững. Saikyou cũng chết lặng bên cạnh tôi, mặt tái nhợt. ``\dots Nii-chan\dots'' giọng em gái tôi lạc đi. ``Anh nhìn xem\dots\ đó có phải\dots\ Touka-san? Rồi Suzune-san, Saki-san\dots\ rồi Inari-san\dots\ và cả Hanabi-san nữa?''

Tôi giật thót, tim lạnh buốt khi nhận ra những gương mặt quen thuộc đang ngồi ngay hàng thẳng lối trong lớp. Một người có bờm tai cáo, một người có tóc hai búi cao, một người tóc xanh nước biển, một người tóc đuôi ngựa buộc lệch, và một người có tóc cột hai bên thấp.

Không thể nào nhầm lẫn đi đâu được. Họ chính là những người đã từng chúng tôi sát cánh bên nhau trong suốt những ngày đầu tiên mà hai anh em lần đầu nhập trường.

``Chuyện này\dots\ không thể nào\dots'' tôi thì thầm, nửa như phủ nhận, nửa như tuyệt vọng. Những người bạn học đã tử nạn sau trận động đất đó, giờ họ đã hóa thành những con rối vô hồn không còn lương tri, và lúc này đang tụ họp lại ở lớp học này.

Sora lại nhìn vào trong, ánh mắt nghiêm trọng khác hẳn lúc nãy. ``\dots Tớ cũng nhận ra vài người tôi từng gặp\dots\ ở những lần trước. Vậy ra\dots\ chúng cũng xuất hiện ở đây.''

``Những lần trước?'' Tôi dám cá là cô ấy đã từng gặp những người khác nữa trước cả tôi, cũng từng trải qua những câu chuyện vừa đau lòng vừa dị hợm này, tới mức mà cảm xúc đã chai sạn đi nhiều.

Nhưng rốt cuộc thì\dots\ Sora là người như thế nào vậy? Tôi vẫn chưa có dịp để hỏi cô ấy, vì cô toàn xuất hiện lúc chúng tôi đang gặp ở những tình huống nước sôi lửa bỏng, tới độ không có thời gian để tìm hiểu nhau.

Tôi định ghé sát hơn\dots

\img{e1-5d}

Ngay lập tức, tất cả bọn họ trong lớp đồng loạt quay đầu về phía chúng tôi. Cùng một lúc. Đôi mắt đen tuyền, xoáy xoáy như bị ai đó điều khiển, khuôn mặt cứng đờ ra, không hề biểu lộ cảm xúc gì. Chỉ có những cặp mắt nhìn về phía chúng tôi, như đang găm chặt bên trong chúng tôi.

``Cái quái--''

*XOẢNG!!*

Ô cửa kính trước mặt tôi bỗng vỡ tan, một vài mảnh sượt qua cánh tay tôi khiến tôi rên nhẹ. Bóng tối ập vào trong lớp như chất lỏng đen đặc. Tiếng cười bật ra, cao, chói, kéo dài, nhọn hoắt xuyên thẳng vào tai, như những kẻ đã hoàn toàn mất trí.

\img{e1-5e}

Tôi nhận ra thứ âm thanh này. Ở những chiều không gian trước đó, và cả trong giấc mơ\dots\ nó luôn đồng nghĩa với điều tồi tệ hơn sắp xảy ra. Tệ hơn rất nhiều so với những gì mà tôi tưởng.

Không chần chừ gì thêm, tôi gầm lên: ``\textbf{Chạy!!}'' kéo mạnh tay Saikyou, đồng thời lao về phía cánh cửa trượt ngay cuối hành lang phía trước mặt.

Sora và tôi bật tung nó ra, rồi nhanh chóng đóng sầm lại.

Tiếng *RẦM!* dội vào tai tôi, tiếp đó là một khoảng tĩnh lặng nghẹt thở. Tim tôi vẫn còn đập loạn xạ sau cảnh tượng kinh hoàng kia, khi mà nó đã thực sự xảy ra, chứ không đơn thuần chỉ trong cơn ác mộng nữa.

Saikyou thở hổn hển, tay chống lên cánh cửa trượt, rồi mắt cô mở lớn khi nhìn thấy vệt máu loang đỏ từ vết rạch do mảnh kính vỡ gây ra.

``\textbf{N-Nii-chan! Tay anh\dots\ máu\dots! Máu chảy rồi kìa!}'' em ấy lắp bắp, khuôn mặt lo âu đến mức gần như sắp khóc.

Tôi cúi xuống xem, thấy đúng là cánh tay mình bị xước, máu rịn ra thành vệt. Tôi thở hắt, cố làm ra vẻ bình tĩnh: ``Chỉ là vết xước nhỏ thôi. Không sao đâu, Saikyou. Đừng lo quá.''

Sora lúc này đã lục lọi trong chiếc túi bên hông. Cô lấy ra một gói bông và băng gạc, điềm nhiên như thể việc này đã quá quen thuộc với mình. ``Để tớ xử lý cho.''

``Ờm\dots\ cậu lúc nào cũng mang theo mấy thứ này à?'' tôi hơi ngạc nhiên, trong khi cô đã nhanh nhẹn lau sạch vết máu và quấn băng quanh cánh tay tôi bằng động tác gọn gàng và dứt khoát.

``Tớ cũng từng bị vài ba lần xây xát như thế rồi, nên tớ hay mang chúng bên mình thôi.''

Saikyou cuối cùng cũng thở phào, đôi mắt rạng lên chút nhẹ nhõm khi cánh tay tôi được băng bó cẩn thận. ``May quá\dots\ Cảm ơn cậu nha, Sora-san.''

Tôi bật cười trừ, nhìn em gái: ``Em đúng là làm quá lên rồi đấy. Chút máu thôi mà, đâu cần phải căng thẳng thế\dots''

``Làm quá á!?'' cô em gái bực dọc, trừng mắt nhìn tôi. Tôi nói gì đó làm phật lòng em ấy sao?

Để rồi, em lại run giọng, đôi mắt chực trào nước: ``Em chỉ\dots\ em chỉ không muốn Nii-chan gặp chuyện gì nữa. Anh là anh trai song sinh của em, em không để anh xảy ra bất trắc đâu. Tuyệt đối không!''

Câu nói ấy khiến tôi thoáng khựng lại. Có chút gì đó vừa ấm áp vừa nặng nề trong tim. Tôi định nói gì đó để xoa dịu thì chợt nghe Sora bật cười khẽ.

``Cô cười cái gì chứ\dots'' tôi gãi đầu, ngượng nghịu.

``Hai anh em cậu thật sự giống nhau thật. Rõ ràng đang sợ chết đi được mà cứ phải mạnh miệng che chở cho nhau đấy,'' cô đáp lại, ánh mắt dịu lại một chút. ``Nhưng mà cũng nhờ hai người mà bầu không khí bớt đáng sợ một chút đó.''

``Thật ra thì\dots\ em cũng sợ lắm,'' Saikyou khẽ cười, ``nhưng có anh ở đây nên em thấy an tâm hơn.''

``Ừm, anh cũng vậy,'' tôi gãi đầu. ``Có em và Sora-san ở đây, anh cũng thấy bớt căng thẳng.''

``Mọi người đừng lo,'' Sora mỉm cười nhẹ. ``Chúng ta sẽ cùng tìm được lối ra thôi.''

``Nhưng mà\dots\ Sora-san này,'' em ngập ngừng. ``Chị có thấy nơi này giống trường cũ không?''

``Ừm, có chút quen quen,'' Sora gật đầu. ``Nhưng khác xa so với những gì tớ nhớ.''

Chúng tôi trao đổi thêm vài câu vu vơ, cố níu lại chút cảm giác bình yên. Nhưng sự nhẹ nhõm ấy, đáng tiếc, không kéo dài lâu.

Ngay khi chúng tôi chuẩn bị tiến bước, một giọng nói khô khốc bỗng vang lên, không phải của bất kỳ ai trong ba người chúng tôi\dots

\dots một giọng có âm điệu tương đồng với Sora-san.

``Cậu đang làm gì ở đây vậy? Nhanh nào, ta cùng về thôi chứ.''

Chúng tôi cùng ngẩng lên.

Ở ngay trước mặt, dưới ánh sáng nhập nhoạng, một cô gái hiện ra. Ngoại hình giống hệt Sora, nhưng đôi mắt xanh sâu thẳm và trên người là bộ đồng phục xanh sẫm cũ kỹ, trông như từ thế kỷ trước. Vẻ mặt lạnh băng, điềm tĩnh đến mức bất an.

Tim tôi lạnh toát. Tôi và Saikyou nhận ra ngay gương mặt ấy. Chính là cô gái chúng tôi đã thấy ở hành lang sau trận động đất kinh hoàng\dots\ trước cả khi gặp được Sora.

Không ai khác, ngoài Misora Shino ra.

Em gái tôi lập tức phản ứng dữ dội, giọng run nhưng kiên quyết: ``\dots Chuyện quái gì đang xảy ra vậy!? Giải thích cho bọn này đi!''

Tôi bước lên một bước, chất vấn thẳng: ``Với cả, tại sao cô lại gọi bọn tôi là `học sinh chuyển trường'? Tại sao chúng tôi lại ở đây?''

Shino im lặng một thoáng. Ánh mắt cô như xuyên qua chúng tôi, nhìn đâu đó xa xăm.

Rồi cô khẽ cất giọng, nhẹ tênh đến lạnh sống lưng: ``Nhanh nào, ta cùng về thôi chứ?''

Tôi nghiến răng: ``Cô đang nói cái gì vậy? `Đi về'? Ý cô là thế nào?''

Saikyou siết chặt nắm tay, hét lớn như trút cả nỗi sợ: ``Bọn tôi sẽ không đi đâu cả cho tới khi cô giải thích tường tận mọi thứ!''

Không khí lập tức đặc quánh lại. Tôi cảm giác từng hơi thở trở nên nặng nề, dính chặt trong lồng ngực. Cả tôi và Saikyou đều quyết liệt. Chúng tôi không còn lùi nữa.

Shino run run người, khuôn mặt thoáng méo mó. Rồi đột ngột tối sầm lại, đôi tay ôm chặt lấy đầu, tỏ vẻ như rất tuyệt vọng. ``Tại sao chứ!? Tại sao cậu lại không trở về!? Cậu phải quay trở về\dots''

Trong con mắt của ả ta, chúng tôi trông không khác gì những đứa trẻ cứng đầu không chịu nghe lời cha mẹ, không chịu làm theo những gì mà người lớn nói.

Và rồi, nước tràn bờ đê. Cái gì cũng có giới hạn của nó. Sự kiên nhẫn của Shino cũng vậy.

Nhưng nếu ả ta chỉ cứ lặp đi lặp lại câu ``về nhà đi''và ``học sinh mới'' (mặc dù thực tế chúng tôi đúng là những học sinh mới chuyển tới) mà không có bối cảnh đầu đuôi rõ ràng, tôi không thể nào chấp nhận được.

Đây là lúc phải ba mặt một lời. Nhưng rồi\dots

*RẮC!*

Một tiếng động khô khốc vọng cả hành lang, như thứ gì đó nứt vỡ.

Toàn thân ả ta rung lên bần bật, như thể tâm trí của ả đã hoàn toàn sụp đổ.

Shino ngẩng phắt mặt lên, đôi mắt mở to hết cỡ, giọng hét vang như xé toạc không gian: ``\textbf{Cậu phải quay về!! Ngay bây giờ!!}''

\img{e1-5f}

Chúng tôi hoảng loạn lùi dần về phía sau. Shino từng bước tiến lên, gương mặt vặn vẹo như một kẻ trốn khỏi viện tâm thần.

``Shino\dots\ Tại sao cậu lại trở nên như vậy?'' Sora thì thào, rồi nhanh chóng đứng chắn giữa chúng tôi và Shino. ``Hai người chạy trước đi!''

``Nhưng mà\dots''

``\textbf{Đừng lo cho tớ! Mau chạy!!}'' cô quát, giọng căng như dây đàn.

Không còn kịp nghĩ, tôi kéo tay Saikyou, cả hai anh em tôi đồng loạt quay đầu bỏ chạy, với Sora-san theo sau.

Tiếng bước chân dội rầm rập trên sàn gỗ, xen lẫn tiếng hét loạn nhịp của Shino đuổi sát phía sau, xen lẫn với những lời nói ma mị không kém:

\begin{center}
  \uline{Cậu đã chạy trốn. Chạy trốn khỏi tội lỗi của mình.}
\end{center}

\begin{center}
  \uline{Cậu đang cố trở thành cái gì? Đã kết thúc rồi\dots\ Đừng có chạy trốn nữa\dots}
\end{center}

Tôi vừa chạy, vừa lia ánh sáng đèn pin điện thoại loạng choạng tìm lối thoát. Những cánh cửa trượt hiện ra liên tiếp, chúng tôi lao đến, kéo mạnh, rồi lại phải bỏ qua vì chỉ toàn phòng tối đặc hoặc ngõ cụt.

Hành lang xung quanh dường như thay đổi không ngừng: có lúc kéo dài vô tận, có lúc ngoằn ngoèo như mê cung. Trần nhà nứt toác, cửa sổ vỡ tung, mảnh gỗ rơi loảng xoảng. Tất cả đều nhuộm trong bóng đen dày đặc, rực sáng bởi ánh đèn pin chập chờn trong tay chúng tôi.

``\textbf{Đừng quay lại! Cứ chạy thẳng!}'' Sora-san hét, giọng gấp gáp. ``Cánh cửa tiếp theo\dots\ Phải có lối thoát gì chứ!''

``Tay anh vẫn ổn chứ, Nii-chan!?'' Saikyou vừa chạy vừa ngoái nhìn tôi, hoảng hốt.

``Giờ không phải lúc để hỏi về cánh tay! Cứ chạy đi!'' tôi gằn giọng, kéo mạnh tay em gái.

Những câu nói vô hồn khàn đặc và tiếng gào thét của Shino vang vọng phía sau, dội theo từng bước chân chúng tôi. Lồng ngực tôi như muốn vỡ tung.

``Hộc hộc\dots''

``Hộc hộc\dots''

Những hành lang kéo dài, hết chỗ này đến chỗ kia, những cánh cửa trượt mà chúng tôi không thể nào đếm nổi cứ thế trôi dần về phía sau trong tiếng gào rú của cô gái tóc nâu mắt xanh kia.

Rồi bất ngờ, một cánh cửa hiện ra ngay phía cuối hành lang. Nó hoàn toàn khác những cánh cửa gỗ cũ kỹ mà chúng tôi đã thấy từ trước. Đây là một cánh cửa trượt màu trắng, bằng kim loại, sáng lạnh, trông rõ ràng thuộc về thời hiện đại.

Tôi thở dốc, mắt mở lớn. ``Có cánh cửa khác thường kìa!''

Tôi không kịp nghĩ nhiều. Chúng tôi nhào tới, kéo mạnh tay nắm. Cánh cửa trượt bằng kim loại trắng phát ra tiếng két khô khốc, rồi từ từ mở ra, để lộ một khoảng không sáng lóa sau lưng nó.

Ngay lập tức, trước mắt tôi hiện ra một khung cảnh\dots\ quen thuộc. Hành lang này\dots\ tôi đã thấy nó rồi. Không phải ở trong giấc mơ, mà là vào cái ngày cả lớp tôi rúng động vì trận động đất. Chính con đường mà tôi và Saikyou lồm cồm rời đi trong lớp 1-C.

Và cũng là nơi\dots\ mà tôi đã nhìn thấy các bạn học đã tử nạn trong trận động đất đó.

``Chúng ta\dots\ quay lại rồi sao?'' tôi lẩm bẩm, thở hổn hển với cảnh tượng quen thuộc.

Cả ba đổ rạp xuống nền, thở dốc như vừa hoàn thành một cuộc chạy marathon. Mồ hôi túa ra trên trán, tay tôi vẫn còn run bần bật khi siết chặt điện thoại. Phía sau, mọi thứ đã im lìm, không còn bóng dáng Shino đâu nữa.

Trước mặt chúng tôi, một cánh cửa trượt màu trắng khác đang mở sẵn. Khác với cánh vừa rồi, từ khe hở phát ra một thứ ánh sáng trắng trong, nhẹ nhõm đến lạ. Chúng tôi không nhìn thấy được bên kia là gì, nhưng nhìn theo ánh sáng kia, tôi có thể dám chắc đó là cách để chúng tôi thoát khỏi đây.

Một sự trùng hợp tới kỳ lạ.

Sora bước lên trước, mắt ánh lên tia hy vọng. ``Đó\dots\ chính là lối ra. Chúng ta có thể thoát khỏi nơi này! Ta có thể quay trở về cuộc sống bình dị rồi!''

Lời cô vang vọng trong hành lang tĩnh mịch, nghe như một lời hứa hẹn, một sự cứu rỗi sau tất cả những gì vừa trải qua.

Tôi và Saikyou nhìn nhau, rồi lặng lẽ đưa mắt trở lại cánh cửa ấy. Ánh sáng trắng ngoài kia hệt như một lối thoát sau cơn ác mộng, nhưng chẳng hiểu sao, trong tôi lại dấy lên một nỗi nghi ngờ mơ hồ.

Sora vẫn đứng trước mặt chúng tôi, bàn tay khẽ giơ ra, như thể đang chờ đợi cả hai bước tới. ``Chúng ta chỉ cần đi qua thôi\dots\ rồi mọi chuyện sẽ bình thường trở lại,'' cô nhấn mạnh một lần nữa, ánh mắt dịu dàng đến mức gần như thuyết phục.

Nhưng chính lúc đó, hình ảnh Hanabi-san, Inari-san, rồi cả Suzune-san và Saki-san\dots\ họ cứ lướt qua tâm trí tôi. Những nụ cười méo mó, những ánh mắt vô hồn trong lớp học vừa rồi\dots\ không thể nào là ``bình thường.''

Tôi siết chặt tay, khẽ hít một hơi. Tôi thực sự không muốn đập tan niềm hy vọng nhỏ nhoi này, nhưng\dots\ có lẽ tôi cũng cần phải thực tế hơn một chút.

``Sora-sam\dots\ nếu thật sự có một `cuộc sống bình thường' chờ ở đó\dots'' tôi gằn giọng, ngập ngừng nhìn Saikyou, ``\dots thì tại sao Hanabi-san và Inari-san lại trở thành\dots\ thứ mà chúng tôi vừa thấy?''

Saikyou tiếp lời ngay, quả quyết hơn tôi: ``Phải đó. Nii-chan nói đúng. Bọn họ chết trong trận động đất, bọn tôi đã thấy tận mắt. Vậy thì cái gì đang chờ ở phía bên kia? Bạn bè của chúng tôi\dots\ hay chỉ là những cái xác vô hồn cứ lặp đi lặp lại một câu?''

Sora-san mở miệng, nhưng không thốt nên lời. Nụ cười nhẹ nhàng ban nãy như bị rút sạch khỏi gương mặt cô. Trong mắt tôi, cô ấy bỗng khựng lại, đứng im, giống như vừa có một sợi dây vô hình siết chặt cổ họng. Cô không phủ nhận, nhưng cũng không dám gật đầu.

Khoảnh khắc ấy, tôi thấy rõ rằng linh cảm của mình - từ giấc mơ định mệnh kia - đã đúng: có điều gì đó mà Sora-san giấu chúng tôi. Có thể là bí mật, có thể là một sự thật tàn nhẫn đến mức cô không đủ can đảm nói ra. Hoặc đúng như Shino nói, cô đang trốn chạy khỏi một sự thật gì đó mà cô không thể chấp nhận.

Tôi cắn môi, một phần muốn tra hỏi thêm, một phần lại thấy thương hại. Saikyou thì siết chặt nắm tay, rõ ràng giận dữ và sợ hãi đan xen. Chúng tôi muốn tin cô ấy, nhưng\dots\ lý trí lại thì thầm rằng chẳng có thứ ``bình thường'' nào tồn tại ở phía bên kia.

Trong dòng suy nghĩ ấy, một giọng nói khô khốc và lạnh lẽo vang lên phía sau lưng làm cả ba chúng tôi giật mình.

``Cậu cứ như vậy mà bỏ chạy sao?''

\img{e1-5g}

Tôi giật thót, quay đầu lại. Bóng dáng đó, khuôn mặt đó\dots\ tôi không thể nào nhầm được. Đôi mắt xanh lạnh băng, mái tóc nâu xõa xuống, gương mặt vô cảm hệt như lần chúng tôi chạm trán trong bóng tối sau trận động đất.

Misora Shino.

Ả ta đã đuổi kịp.

``Sora\dots'' Shino cất tiếng, giọng đều đều mà như ghì chặt, ``cậu sẽ cứ mãi trốn tránh ư? Bao lâu nữa thì thôi? Cậu có thể giả vờ rằng tất cả vẫn ổn, nhưng thực tại thì chẳng bao giờ thay đổi.''

Bầu không khí xung quanh bốn người chúng tôi bỗng chốc chùng xuống. Hai con người, hai bản sao gần giống nhau, nhưng tính cách hoàn toàn trái ngược, giờ đây đang đứng đối diện với nhau, với tâm thế cũng ngược nhau. Một bên là sự dè chừng và có phần sợ sệt, và một bên thì trông như nắm rõ mọi thứ trong lòng bàn tay.

Tôi và Saikyou đứng bên ngoài, vô tình đứng làm ``khán giả bất đắc dĩ.'' Dường như ả ta không đả động gì chúng tôi - hoặc ít nhất, tôi nghĩ vậy.

Sora quay lại, ánh mắt lay động, nhưng vẫn gắng giữ sự bình tĩnh: ``Không. Tớ không chối chạy. Tớ\dots\ tớ chỉ muốn bảo vệ họ. Muốn họ có cơ hội được sống tiếp.''

Tôi lặng im. Trong khoảnh khắc đó, tôi nhớ lại giấc mơ đêm nhập học. Ở đó, tôi đã thoáng thấy cảnh này: hai người họ đối đáp nhau, từng lời như mũi dao đâm xuyên qua nhau. Chẳng lẽ giấc mơ ấy là điềm báo cho thứ đang xảy ra ngay trước mắt tôi?

``Bảo vệ?'' Shino nhếch môi cười nhạt, ánh mắt lóe lên thứ gì đó cay nghiệt. ``Cậu gọi việc dẫn họ vào vòng lặp bất tận này là bảo vệ sao? Cậu chỉ đang trốn tránh khỏi thực tại, rồi tuyệt vọng níu kéo lại những mảnh vụn đã vỡ nát mà thôi.''

``Cậu cũng có hơn gì tớ đâu?'' Sora bật lại, giọng run nhưng dứt khoát. ``Cậucho rằng ép họ quay về với cái chết mới là đúng sao!? Chẳng lẽ chỉ có nỗi đau và sự mất mát mới được phép tồn tại?''

Hai người đứng đối diện nhau, căng thẳng đến mức từng hơi thở cũng nặng nề. Tôi và Saikyou chỉ biết đứng chen giữa, đôi mắt tôi dán vào họ như bị thôi miên. Cuộc tranh cãi này\dots\ vốn không thuộc về chúng tôi, nhưng lại vô tình cuốn chúng tôi vào như một dòng chảy không thể cưỡng.

Trong cuộc nói chuyện giữa hai cô gái có ngoại hình tương đồng, một giọng nói nhỏ nhẹ chợt cất lên, áo tôi như bị giật nhẹ: ``Nii-chan\dots\ Nhìn kìa.''

Saikyou run run chỉ tay xuống sàn. Ngay bên cạnh chân cô, một bông hoa đỏ thẫm mọc lên từ khe nứt của gạch lát.

Một bông hoa bỉ ngạn đỏ.

\img{e1-5i}

Thoạt nhìn, nó đẹp đến nghẹt thở. Những cánh hoa mảnh khảnh như tơ, tỏa ra thành chùm đỏ rực, như đốm lửa trong bóng tối đặc quánh. Không ngạc nhiên mà Saki-san là người đã cất công trồng chúng bên ban công, quý trọng chúng như những người con, đù rằng tôi thấy hầu như không ai trồng loài hoa này.

Nhưng trong lòng tôi chợt thắt lại. Tôi nhớ đến những gì từng nghe, từng đọc: bông hoa của chia ly, của sự tiễn biệt, của những linh hồn bơ vơ. Một vẻ đẹp đẫm buồn, gắn liền với những khúc ca tang tóc.

Tôi từng nhìn thấy bông hoa này được đặt trong một cái lọ, và được đặt trên một bàn học. Ở Nhật Bản, việc đặt lọ hoa trên bàn của một học sinh thường sẽ mang hai ý nghĩa phản nghịch nhau.

Có thể là họ đặt ở đó để tưởng niệm học sinh ấy đã qua đời.

Hoặc tệ hơn\dots\ là một cách để nguyền rủa người đó xuống dưới mồ. Một cách bắt nạt thầm lặng nhưng đủ khiến đối phương rợn người.

Saikyou cúi xuống, khẽ nâng bông hoa ấy trên tay, rồi khẽ cất tiếng: ``H-Hai người\dots'' Ngay lập tức, Shino và Sora đồng loạt im bặt.

Phản ứng của họ trái ngược hoàn toàn.

Shino trầm ngâm, ánh mắt tối lại như vừa nhìn thấy một mảnh ký ức cũ ùa về, một thứ gì đó ả đã chứng kiến không chỉ một lần. Tôi thấy rõ sự quen thuộc trong ánh nhìn ấy, như thể bông hoa là dấu mốc gắn liền với câu chuyện của riêng mình.

Trong khi đó, Sora thoáng lùi lại, đôi môi mím chặt, rõ ràng có chút sợ sệt. Có vẻ cô ấy cũng biết rất rõ bông hoa ấy biểu trưng cho cái gì.

Cả hành lang chìm vào tĩnh lặng, chỉ còn bông hoa đỏ rực kia như một ngọn lửa chập chờn soi sáng khoảng cách mong manh giữa họ.

Cô gái có đôi mắt xanh lặng đi. Một thoáng trầm mặc, cô khẽ nhắm mắt, rồi mở ra với ánh nhìn chất chứa thứ gì đó\dots\ vừa u ám, vừa rùng rợn và có phần bi thương.

``\dots Hoa bỉ ngạn,'' Shino cất tiếng, giọng khàn khàn như gió lùa qua vách gỗ mục. ``Nó luôn xuất hiện\dots\ mỗi lần những chuyện tồi tệ nhất xảy đến trong đời tôi.''

Câu nói ấy khiến tôi thoáng lạnh gáy. Tôi ngập ngừng hỏi: ``Ý cô là\dots\ quá khứ của cô gắn liền với thứ này sao?''

Saikyou liếc sang, khẽ gật đầu tiếp lời tôi: ``Nếu vậy, nó rốt cuộc\dots\ có ý nghĩa gì? Tại sao lại gắn với cô đến mức này?''

Ả bật cười nhạt, nụ cười méo mó không mang chút ấm áp nào: ``Các cậu sẽ chẳng hiểu đâu, nếu tôi chỉ kể suông. Chữ nghĩa hay từ ngữ sẽ chỉ làm nghèo nàn đi sự thật. Có những điều\dots\ chỉ khi tự mình trải qua, mới có thể thấu hết được.''

Tôi sững người. Trong giọng điệu ấy có một tầng nghĩa khác, ẩn sâu hơn. Không phải chỉ kể chuyện đơn thuần, ả ta\dots\ đang ép chúng tôi bước vào nó.

Tôi không nghĩ Shino sẽ sẵn sàng mở lòng với chúng tôi. Hẳn Shino phải tìm thấy một điểm nào đó tương đồng với ả thì mới có thể nói như vậy. Nói cách khác, chỉ những người cùng nằm trong hoàn cảnh ấy mới hiểu được vấn đề.

\dots

\dots

``Học sinh mới chuyển trường.'' Đúng, là thứ gắn vào lời nguyền của ngôi trường.

Cụm từ này cứ văng vẳng hết từ lời truyền tai trong trường, rồi đến những mẩu báo, những tờ ghi chú mà tôi và Saikyou đã đọc qua trong giấc mơ.

Cụm từ ấy\dots\ hẳn phải gắn với một sự kiện quỷ quái nào đó nên mới thành ra như vậy.

Tôi ngước nhìn về cô em gái, người vẫn đang nắm chặt bông hoa bỉ ngạn. Em ấy không nói gì, chỉ giương đôi mắt xanh về phía tôi, và khẽ gật đầu, như thể hai chúng tôi có cùng chung một suy nghĩ.

Ta phải tìm ra đầu đuôi của mọi chuyện. Có như thế, chúng tôi mới có thể hiểu được tại sao lại có lời nguyền kia.

Shino đưa mắt nhìn chúng tôi, đôi đồng tử xanh ánh lên ánh sáng khác thường. ``Để các cậu hiểu\dots\ tôi sẽ cho hai người thấy khoảng thời gian mà tôi còn sống.''

Tôi và Saikyou đồng loạt khựng lại. Tim tôi như bị ai đó bóp chặt. ``Còn sống''\dots? Tức là\dots

``\dots cô đã chết rồi ư?'' cả hai anh em tôi đồng thanh, mắt trố ra như không tin vào sự thật. Tôi đánh mắt về phía Sora, vẫn đang nắm chặt vào vạt váy, im lìm từ nãy tới giờ.

Cô ấy không nói một câu nào. Tôi biết, khuôn mặt đó là một vẻ cam chịu, trông như không muốn dính dáng vào chuyện này, nhưng\dots\ không phải là Sora không đồng tình với ả. Cô ấy chỉ muốn rời khỏi đây, càng sớm càng tốt.

Shino không nói gì thêm. Chỉ khẽ nhấc cằm, nhìn thẳng vào chúng tôi. Ả hướng bàn tay về phía trước, cất giọng cảnh báo với vẻ nghiêm nghị: ``Hai anh em cậu\dots\ thật sự sẽ phải cẩn thận đấy. Nếu không cẩn thận, các cậu có thể sẽ bỏ mạng tại đó, \emph{Hikawa-san}.''

\img{e1-5h}

Ả sau đó quay lưng đi, để lại một bầu không khí đặc quánh.

``Hikawa-san?''

Tôi không điếc. Từng âm được phát ra rõ ràng, dứt khoát. Tôi và Saikyou thoáng giật mình, nhưng chẳng kịp hỏi lại vì sao cô biết. Bóng dáng của ả đã nhạt dần trong khoảng tối của phía cuối hành lang từ khi nào.

Để rồi sau đó\dots

*Ù Ù Ù\dots!*

Ngay khi đôi mắt chúng tôi mất dấu cô ta, mặt đất bất chợt rung chuyển. Cơn chấn động lan nhanh như cuộn sóng ngầm, từng tấm trần rạn nứt, bụi trần rơi rào rào, đi kèm với những tấm ốp lát phía trên.

``Động đất nữa ư\dots!?'' tôi nghiến răng, chồm tới kéo tay Saikyou. Nhưng chúng tôi đã không kịp phản ứng.

*RẮC!*

Một vết nứt khổng lồ toác ra ngay dưới chân, kéo chúng tôi thẳng xuống một khoảng không đen kịt.

``\textbf{Nii-chan!!!}'' em tôi hét thất thanh, còn tôi chỉ kịp cảm nhận cơ thể rơi tự do.

``\textbf{Saikyou! Bám chặt vào anh!}'' tôi hét lớn đáp lại, quờ quạng trong ánh sáng đang dần tắt lịm.

Trong khoảnh khắc, tôi nghe thấy tiếng Sora từ phía trên vọng xuống, gào xé như tuyệt vọng: ``\textbf{Kyuen-san! Saikyou-san!! Đừng đi!!!}''

Rồi tất cả biến mất, chỉ còn lại tiếng ù ù của hư không\dots\ và khoảng không đen sâu hun hút kéo chúng tôi vào.

Bóng tối\dots\ nuốt chửng tất cả. Một lần nữa. Bỏ lại tất cả mọi thứ đang dần sụp đổ trên kia.

Gió lạnh tạt thẳng vào mặt, nhưng trong tâm trí, tôi chẳng còn nghĩ được gì ngoài những mảnh ghép hỗn loạn.

Hoa bỉ ngạn đỏ\dots\ lời Shino buông ra trong thứ giọng trầm nặng, nửa như một lời trăn trối, nửa như một lời nguyền. Quá khứ của cô, gắn liền với cái chết và sự chia ly, rồi cả lời cảnh báo rằng chúng tôi có thể sẽ bỏ mạng. Tôi nghe đi nghe lại trong đầu, mà chẳng biết điều gì mới là thật.

Nếu quả thực có một ``lời nguyền'' ám lên ngôi trường này, nó hẳn không dừng lại ở những kẻ đã chết như Hanabi-san hay Inari-san. Nó còn vươn tới cả tôi và Saikyou, những kẻ mới chỉ vừa đặt chân đến.

Nhưng\dots\ thứ khiến tôi bứt rứt hơn cả không phải là hoa bỉ ngạn, cũng chẳng phải là bóng dáng Shino đang khuất dần trong ký ức. Mà chính là cái tên cô gọi trước khi rời đi.

``Hikawa-san.''

Âm thanh đó vẫn vang lên rất rõ ràng, như khắc sâu vào thính giác. Một cái tên xa lạ, chẳng hề dính dáng đến chúng tôi. Hay\dots\ nó vốn dĩ chính là một phần của chúng tôi, mà vì lý do nào đó, cả hai đã quên mất?

Tôi không tìm thấy câu trả lời. Chỉ có sự trống rỗng lan dần trong lồng ngực. Tôi thoáng liếc sang cô em gái, thấy em cũng đang mở to đôi mắt, không giấu nổi sự băn khoăn giống hệt tôi.

Và rồi, như chẳng còn cách nào khác, chúng tôi để mặc bản thân rơi, trượt dài vào lòng bóng tối.

Không đáy. Không ánh sáng. Chỉ còn tiếng tim đập, và ý nghĩ chồng chéo về những gì mà cô gái mắt xanh kia nói, dần nuốt trọn tôi và em gái, khi cả hai chìm sâu vào màn đêm đen kịt không lối thoát.

\end{document}